% Primitive: random-variables — authored from
% probability/random_variables_manim.py, its README concepts table (six
% scenes) and plan 007's verified anchors (main report + addendum). Closes
% the road's oldest promise: the counting chapter's C(n,k) returns as a
% count of sorted-square columns. Problem answers are computed by
% answers/random_variables.py (the solve-gate anchor); every exact value
% there is Fraction arithmetic, floats only where a decimal display is the
% claim.
\chapter{Random variables}

\begin{narrative}
So far the road weighs \emph{events}: regions of a square, shares of
area. But the questions it is heading toward are about \emph{numbers}
read off outcomes — how many heads, which symbol, how large a score. A
random variable attaches a number to every outcome; this chapter builds
the attachment, watches it sort into a distribution, and finds the one
number to act on.

\section{The stamped square}

A die is usually introduced as a set of six outcomes. This series
opens differently: the die as a \emph{function} — $X(\text{face}) =$
the number shown — ink stamped on the sample space before anything is
rolled~\cite{probability-joe-blitzstein-stat-110}. Then the square the
independence chapter quartered twice returns for its third appearance:
sixteen equal cells, one per four-flip sequence, each stamped with its
head count, $0$ through $4$. The stamp is fixed; the only random
object is where the dart lands; the label is \emph{looked up}, never
generated. Only then, the formula:
\[
  X : \Omega \to \mathbb{R} .
\]
Every measurement attached to a random process is this — a fixed rule
reading a random outcome.

\section{Sort the square}

Slide the stamped cells into columns grouped by their stamp. Area is
conserved — nothing is created, nothing weighed twice — and the bars
that form are the distribution: five \emph{unequal} bars born from
sixteen \emph{equally likely} cells,
\[
  P(X = k) \;=\; \frac{\#\{\omega : X(\omega) = k\}}{16}
  \;=\; \anchor{007.pmf16} .
\]
The unequal weights are visible counts of equal cells — which is the
working refutation of the equiprobability bias, the documented habit
of reading randomness as uniformity that survives instruction
unless the mechanism is
shown~\cite{probability-gauvrit-morsanyi-the-equiprobability-bias-pmc}.

Then the blueprint-versus-house beat: stamp the \emph{same} square
with $Y =$ tails. The sorted columns are identical — same pmf — yet
$X \neq Y$, and $X + Y = 4$ in every single cell. The distribution
forgets which cell was which; the variable remembers. Mistaking one
for the other is the canonical error this scene exists to prevent.

\section{The balance point}

Expectation is defined, not simulated. Hang the die's six equal bars
on a beam and ask where the fulcrum goes: it balances at
$\anchor{007.dieE}$ — which is not a face. ``The point at which the
distribution would balance'' is the picture; a weighted average need
not be attainable~\cite{probability-orloff-bloom-mit-18-05-class-4b-expectation}.

\begin{center}
\begin{tikzpicture}[x=1.15cm, y=0.5cm]
  % SortTheSquare's columns — sixteen equal cells grouped by stamp —
  % with TheBalancePoint's fulcrum under the axis at E[X] = np = 2.
  \foreach [count=\k from 0] \c in {1,4,6,4,1} {
    \foreach \i in {1,...,\c}
      \draw[Muted, fill=CoolInk!25]
        (\k-0.42, \i-0.92) rectangle (\k+0.42, \i-0.08);
    \foreach \i in {1,...,\c}
      \node[color=CoolInk!30!black] at (\k, \i-0.5) {\scriptsize $\k$};
    \node[above, color=CoolInk] at (\k, \c+0.25) {\small $\c/16$};
  }
  \draw[Muted] (-0.7, -0.55) -- (4.7, -0.55);
  \fill[AccentInk] (2, -0.75) -- (1.72, -1.9) -- (2.28, -1.9) -- cycle;
  \node[right, color=AccentInk] at (2.62, -1.35) {\small $E[X] = 2$};
\end{tikzpicture}
\end{center}

For the flips, the two ways of summing are shown equal once, and the
sort is the proof: summing stamps over cells gives $32/16 = 2$, and
summing values times column weights gives $2$ — the sort was a
\emph{regrouping}, which is why
\[
  E[X] \;=\; \sum_x x \cdot P(X = x)
\]
may be computed cell by cell or column by column. And the balance
point belongs to the \emph{measure}, exactly as independence did: put
the independence chapter's biased die on the beam — double weight on
the six — and
the fulcrum moves to $\anchor{007.biasedE}$, again not a face. The
level-three reading is Huygens, 1657, the first printed probability
work: expectation is the fair price of the ticket — the number you
act on, with no long run in
sight~\cite{probability-j-j-o-connor-and-e-f-robertson-christiaan-huygens-mactutor}.

\section{Same outcomes add}

$E[X+Y]$ is one sum over the \emph{same} outcomes, and addition
distributes — so
\[
  E[X + Y] \;=\; E[X] + E[Y]
\]
with independence nowhere in the hypothesis: as Grinstead and Snell
put it, mutual independence of the summands was not
needed~\cite{probability-grinstead-snell-introduction-to-probability-6-1-libretexts}.
The demonstration runs at the dependent extreme: $X$ and $4 - X$ are
as dependent as two variables can be, and their sum is $4$, always.
Then the independence chapter's $6 \times 6$ grid is re-read: the
two-dice sum paints diagonals, distribution
$\anchor{007.twodice}$, and $E = 7$ arrives twice — once by the
diagonals' weighted sum, once as
$\anchor{007.dieE} + \anchor{007.dieE}$. Linearity is the workhorse:
it prices any bundle from its parts.

\section{The binomial columns}

Now the road's oldest promise closes. Count cells per column — the
numerators of $\anchor{007.pmf16}$ — and notice that counting
four-letter H/T words with $k$ heads is exactly the counting
chapter's $\binom{4}{k}$. The coefficient is a \emph{cell count},
never a memorized factor.

Re-cut the square with the cuts at $p = 1/4$. The cuts stay straight
— straight cuts are the square's grammar for independence, and
independence is unchanged; only the cut positions move. The cells go
unequal, but every $k$-head cell keeps the \emph{same} area
$p^k q^{4-k}$, because a product does not care about factor order. So
column $k$ weighs (cell count) $\times$ (one cell's area), and the
pmf assembles with nothing smuggled:
\[
  P(X = k) \;=\; \binom{n}{k}\, p^k (1-p)^{n-k} ,
\]
at $p = 1/4$ concretely $\anchor{007.pmf256}$, each one-head cell
exactly $\anchor{007.onecell}$ with $\binom{4}{1} = 4$ such cells.
Expectation costs no combinatorics at all: stamp each trial's
indicator and add four expectations of $p$ — $E = np$, here $1$. The
model's conditions are named, not assumed~\cite{probability-illowsky-dean-openstax-introductory-statistics-2e-4-3}:
fixed $n$, two outcomes, constant $p$, independent trials — and the
independence chapter's aces already marked the boundary: no
replacement means no binomial. One more door opens here: zero
successes in $n$ trials of chance $1/n$ crowds
$1/e \approx \anchor{007.oneovere}$ from below — the previous
chapter's ceiling, mirrored in a shrinking power.

\section{Proportions converge}

The closer turns the swamping intuition into exact numbers — exact
binomial sums, no new machinery. Within $\pm 5\%$ of half the flips:
the probability \emph{climbs}, $\anchor{007.propband}$ at
$n = 20, 100, 1000$. Within $\pm 5$ \emph{heads} of half: it
\emph{falls}, $\anchor{007.countband}$ at $n = 100, 1000, 10000$. The
two $n = 100$ entries are the same band of outcomes, $45$ to $55$ —
one number telling two stories. Proportions converge while counts
spread: the gambler's fallacy dies by two columns moving in opposite
directions. Chance \emph{swamps}; it does not compensate.

The theorem these tables compute instances of is the weak law of
large numbers — Bernoulli had it proved by about 1689; it reached
print in 1713~\cite{probability-grinstead-snell-ch-8-source,
probability-j-j-o-connor-and-e-f-robertson-jacob-bernoulli-mactutor}
— named here, and promised for the day variance arrives. One small
closing count: average the logarithms chapter's surprisal over the
sixteen equal cells and it comes to exactly $4$ bits — a number
waiting for the information-theory road.

\paragraph{Where this goes.} The next move pins the data and sweeps
the parameter: the likelihood lens that opens the softmax chapter
reads a binomial table along the other axis, and this chapter's
columns are that table. The per-frame distributions the trellis
multiplies are pmfs exactly like these — the gradient chapter
differentiates straight through them. And the weak law returns as a
theorem when spread gets its own machinery: variance and the proof
arrive together.
\end{narrative}

\section*{Practice problems}

\begin{problem}
Three fair flips, $X =$ number of heads. Stamp the eight cells and
sort: what is the pmf of $X$?
\begin{hint}
Eight equal cells; group by stamp and count.
\end{hint}
\begin{solution}
The eight sequences stamp as one $0$, three $1$'s, three $2$'s and
one $3$, so the pmf is $(1, 3, 3, 1)/8$ — and the counts are
$\binom{3}{k}$, the counting chapter's H/T words one flip early.
Four unequal bars from eight equal cells: the equiprobability bias
refuted at $n = 3$.
\end{solution}
\end{problem}

\begin{problem}
Where does the fulcrum balance for a fair die? And for the biased
die with double weight on the six (weights $1,1,1,1,1,2$ over $7$)?
\begin{hint}
$E[X] = \sum x \cdot P(X = x)$, exact fractions throughout.
\end{hint}
\begin{solution}
Fair: $E = (1 + \cdots + 6)/6 = 21/6 = 7/2 = 3.5$ — not a face. Biased:
$E = (1 + 2 + 3 + 4 + 5 + 2 \cdot 6)/7 = 27/7 \approx 3.857$ — also
not a face, and \emph{moved}: the balance point belongs to the
measure, not to the faces.
\end{solution}
\end{problem}

\begin{problem}
Two fair dice, $S$ = their sum. Compute $E[S]$ twice: once from the
distribution of $S$, once with no distribution at all.
\begin{hint}
The $6 \times 6$ grid's diagonals carry the counts
$(1,2,3,4,5,6,5,4,3,2,1)$; linearity needs neither them nor
independence.
\end{hint}
\begin{solution}
By the distribution: the diagonal counts over $36$ give
$E[S] = \sum_{s=2}^{12} s \cdot P(S = s) = 252/36 = 7$. By linearity:
$E[S] = E[X_1] + E[X_2] = 3.5 + 3.5 = 7$, using only that both dice
are fair — the second route never asks whether the dice are
independent, because it does not need to know.
\end{solution}
\end{problem}

\begin{problem}
Assemble the binomial pmf for $n = 4$, $p = 1/4$ the sorted-square
way: cell counts times one cell's area. Then check $E = np$ by the
definition sum.
\begin{hint}
The counts are $1, 4, 6, 4, 1$; a $k$-head cell has area
$(1/4)^k (3/4)^{4-k}$.
\end{hint}
\begin{solution}
Column $k$ weighs $\binom{4}{k} (1/4)^k (3/4)^{4-k}$:
\[
  \frac{81}{256},\ \frac{108}{256},\ \frac{54}{256},\
  \frac{12}{256},\ \frac{1}{256},
\]
summing to $256/256 = 1$. Each one-head cell is
$(1/4)(3/4)^3 = 27/256$, and there are $\binom{4}{1} = 4$ of them —
coefficient as cell count, power as one cell's area. The definition
sum gives
$E = (0 \cdot 81 + 1 \cdot 108 + 2 \cdot 54 + 3 \cdot 12 + 4 \cdot 1)/256
= 256/256 = 1$; indicators agree instantly: $np = 4 \times 1/4 = 1$.
\end{solution}
\end{problem}

\begin{problem}[stretch]
Zero successes in $n$ trials of chance $1/n$: compute
$(1 - 1/10)^{10}$ exactly and $(1 - 1/100)^{100}$ to six places, and
compare both with $1/e$. From which side, and in what order, do they
approach it?
\begin{hint}
$(9/10)^{10}$ terminates as a decimal; the last chapter named the
limit.
\end{hint}
\begin{solution}
$(9/10)^{10} = 3486784401/10^{10} = 0.3486784401$ exactly, and
$(99/100)^{100} \approx 0.366032$, against
$1/e \approx 0.367879$. Both sit \emph{below} $1/e$ and climb toward
it — $0.348678 < 0.366032 < 0.367879$ — the split-year ceiling
mirrored: $(1 + 1/n)^n$ crowds $e$ from below, and
$(1 - 1/n)^n$ crowds $1/e$ the same way. Rare events at rate one-per-
trial-count leave you empty-handed about $37\%$ of the time, almost
regardless of $n$.
\end{solution}
\end{problem}
